\section*{NeurIPS Paper Checklist}

\begin{enumerate}

\item {\bf Claims}
    \item[] Question: Do the main claims made in the abstract and introduction accurately reflect the paper's contributions and scope?
    \item[] Answer: \answerYes{}.
    \item[] Justification: The abstract and introduction claim a conformance-gated local benchmark and explicitly limit the stability claim; Tables~\ref{tab:portable-main-results} and~\ref{tab:portable-decision-table} provide the supporting evidence.

\item {\bf Limitations}
    \item[] Question: Does the paper discuss the limitations of the work performed by the authors?
    \item[] Answer: \answerYes{}.
    \item[] Justification: The manuscript includes a dedicated limitations section describing the single-node setting, bounded S3 corpus, S2 freshness interpretation, and budget-stability limits.

\item {\bf Theory assumptions and proofs}
    \item[] Question: For each theoretical result, does the paper provide the full set of assumptions and a complete proof?
    \item[] Answer: \answerNA{}.
    \item[] Justification: The paper does not present theoretical results; it defines benchmark metrics and decision rules.

\item {\bf Experimental result reproducibility}
    \item[] Question: Does the paper fully disclose all the information needed to reproduce the main experimental results?
    \item[] Answer: \answerYes{}.
    \item[] Justification: The appendix and artifact README list the workflow commands, archive identifier, validation command, runtime metadata, checksums, report artifacts, and local artifact-readiness verifier.

\item {\bf Open access to data and code}
    \item[] Question: Does the paper provide open access to the data and code, with sufficient instructions to faithfully reproduce the main experimental results?
    \item[] Answer: \answerYes{}.
    \item[] Justification: The benchmark code, commands, generated manifests, and archived result bundle are part of the repository prepared with this manuscript; any public release should preserve anonymization at submission time.

\item {\bf Experimental setting/details}
    \item[] Question: Does the paper specify all training and test details necessary to understand the results?
    \item[] Answer: \answerYes{}.
    \item[] Justification: The benchmark-design section specifies engines, datasets, embedding tiers, concurrency pins, budget levels, metrics, decision rules, and conformance gates.

\item {\bf Experiment statistical significance}
    \item[] Question: Does the paper report error bars suitably and correctly defined or other appropriate information about statistical significance?
    \item[] Answer: \answerYes{}.
    \item[] Justification: The main table reports 95\% confidence intervals; targeted S2/S3 audits report paired bootstrap intervals for the highest-risk margins.

\item {\bf Experiments compute resources}
    \item[] Question: For each experiment, does the paper provide sufficient information on compute resources needed to reproduce the experiments?
    \item[] Answer: \answerYes{}.
    \item[] Justification: The scope and reproducibility sections state the single-node commodity-hardware setting and exact budget ladder; the archived run metadata records per-run resource fields.

\item {\bf Code of ethics}
    \item[] Question: Does the research conducted in the paper conform with the NeurIPS Code of Ethics?
    \item[] Answer: \answerYes{}.
    \item[] Justification: The work evaluates retrieval infrastructure using public benchmark artifacts and does not involve human subjects, private data, or model training on sensitive data.

\item {\bf Broader impacts}
    \item[] Question: Does the paper discuss both potential positive societal impacts and negative societal impacts of the work performed?
    \item[] Answer: \answerYes{}.
    \item[] Justification: The limitations and broader-impact section discusses the benefit of transparent deployment evidence and the risk of overgeneralizing local benchmark winners.

\item {\bf Safeguards}
    \item[] Question: Does the paper describe safeguards for responsible release of data or models that have a high risk for misuse?
    \item[] Answer: \answerNA{}.
    \item[] Justification: The paper does not release new generative models or scraped high-risk datasets; it uses existing retrieval benchmark assets.

\item {\bf Licenses for existing assets}
    \item[] Question: Are creators or original owners of assets properly credited and are licenses and terms of use explicitly mentioned and respected?
    \item[] Answer: \answerYes{}.
    \item[] Justification: The manuscript cites the datasets, software systems, and embedding models used; final submission should include any repository-level license manifests alongside the artifacts.

\item {\bf New assets}
    \item[] Question: Are new assets introduced in the paper well documented and is the documentation provided alongside the assets?
    \item[] Answer: \answerYes{}.
    \item[] Justification: \textsc{HotpotQA-portable} is documented through its preprocessing description, manifest, checksums, archived bundle, and Croissant-style metadata included with the reviewer artifact.

\item {\bf Crowdsourcing and research with human subjects}
    \item[] Question: For crowdsourcing experiments and research with human subjects, does the paper include instructions, screenshots, and compensation details?
    \item[] Answer: \answerNA{}.
    \item[] Justification: The paper does not involve crowdsourcing or human-subject experiments.

\item {\bf Institutional review board approvals or equivalent for research with human subjects}
    \item[] Question: Does the paper describe risks to study participants and IRB approvals or equivalent?
    \item[] Answer: \answerNA{}.
    \item[] Justification: The paper does not involve human-subject research.

\item {\bf Declaration of LLM usage}
    \item[] Question: Does the paper describe the usage of LLMs if it is an important, original, or non-standard component of the core methods?
    \item[] Answer: \answerNA{}.
    \item[] Justification: LLMs are not an important, original, or non-standard component of the benchmark method; the benchmark evaluates retrieval infrastructure and local embedding models.

\end{enumerate}
